%% -*- coding:utf-8 -*-

\section*{Foreword of the sixth edition}


% Added Section [Section~4.2.4.1] to references to ATB extraction and Ross67 udc.tex
% This went to the rerevised edition5

I added and fixed references and added information like page numbers or information on sections to
some references.

% 22.02.23
I changed the adjunct rule (\ref{ex-LFG-adjuncts-rule}) in the chapter on LFG so that adjuncts are attached to VP rather than
V$'$ (p.\,\pageref{ex-LFG-adjuncts-rule}). This is correct for English and for German, since in Berman's approach to German, all
projections combining verbs with arguments are VPs.
% In LFG changed adjunct rule to VP rather V'. This gets the German adjuncts right as well,
% since Berman has all verbal projections be VPs.
%
% Removed reference to Saussure after mentioning competence/performance distinction.
% 21.06.2023
% Thanks to Patrik Austin

% 28.08.2023
% Added page number for SW2011a
% Changed These into Examples in the intro of phrasal.tex since I discuss only some examples of
% phrasal constructions. There are probably more.

% 07.12.2023 fixed reference to SKS2013, which is KSS2013 in 3gb.tex on lowering
%
% Added \citealp[\page 615--616]{CDGRP2013a-u}; as reference to inheritance-based analysis
% of the passive in TAG
%
% 19.12.2023
% GPSG was not context-sensitive but context-free
% since 2018
Section~\ref{sec-recursion-empirical-problems} contained an embarrassing typo. I said that GPSG was designed to be
context-sensitive (p.~463 of the second edition), which did not make any sense in the context of the
following sentence: GPSG was designed to be context-free, but later it was shown that this is not
powerful enough. This is now fixed (p.\,\pageref{page-GPSG-context-free}).


% 05.03.2024
The word in (\ref{ex-categories-with-features}) on p.\,\pageref{ex-categories-with-features} was
changed from \emph{des} to \emph{der}. Thanks to Hannah Schleuper for
pointing out this mistake, which was in the books since the first German edition in 2010. 

%05.10. Antonio zum Grundkurs Reader ein interessantes -> das interessante
% {ex-nounless-np} und {Beispiele-NP-Adjunkte}

I added Chapter~\ref{chapter-LLMs} on Large Language Models and the status of theoretical
linguistcs. This chapter grew out of a discussion note in the Zeitschrift für Sprachwissenschaft. I thank Rui Chaves, Mark Felfe, Hubert Haider, Joachim Jacobs, Tibor Kiss, Alexander Koller, Roland
Schäfer, Oliver Schallert and Remi van Trijp, Giuseppe
Varaschin for comments and discussion on that paper and Konstantin Schulz, Anke Lüdeling, and Martin Klotz for 
discussion. Special thanks go to Christine Mooshammer for raising questions about
\citew{AmbridgeBlything2024a}. This got me started on the topic of LLMs.

%26.04.2025 JS1997a als Quelle für Beispiel mit Fernabhängigkeiten angegeben
% John is Kim now in TAG chapter.
%
% 24.05.2025 added references for lexical rules/metarules in TAG chapter
% fixed layout issue in D-trees in TAG chapter


Figure~\ref{Abbildung-Haugereid} on p.\,\pageref{Abbildung-Haugereid} contained a mistake since the first edition of this book. It is fixed now.

~\medskip

\noindent
Berlin, 2025\hfill Stefan Müller



%      <!-- Local IspellDict: en_US-w_accents -->
